\documentclass[journal]{IAENGtran}
\usepackage{amsmath}
\usepackage{amssymb}
\usepackage{graphicx}
\usepackage{orcidlink}

\hyphenation{op-tical net-works semi-conduc-tor}
\emergencystretch=3em

\newtheorem{Theorem}{Theorem}[section]
\newtheorem{Lemma}[Theorem]{Lemma}
\newtheorem{Proposition}[Theorem]{Proposition}
\newtheorem{Corollary}[Theorem]{Corollary}
\newtheorem{Definition}[Theorem]{Definition}
\newtheorem{Remark}[Theorem]{Remark}

\begin{document}

\title{Limits of Self-Correction in LLMs: When Self-Evaluation Accepts Wrong Answers}

\author{Andrew~Michael~Brilliant,
        \orcidlink{0009-0004-8024-5442}%
\thanks{A.~M.~Brilliant is with Applied Dynamics Research, Sapporo, Japan (e-mail: ab@ad-research.org).}% <-this % stops a space
}

\markboth{International Journal of Computer Science (IJCS)}%
{Brilliant: Limits of Self-Correction in LLMs}

\maketitle

\pagestyle{empty}
\thispagestyle{empty}



\begin{abstract}
We measure whether self-evaluation catches errors. Each task contributes one immutable
candidate, hashed before evaluation and never modified; evaluators judge without
rewriting. For each model family we report the false-acceptance rate
$\kappa=\Pr(A{=}1\mid T{=}0)$, the true-acceptance rate $\tau$, and the
discrimination gap $\Delta=\tau-\kappa$, backed by a shared-blind-spot information
bound $I(T;S\mid G)\leq I(T;Z\mid G)$ and a companion proposition on treating
coupled verdicts as independent. A HumanEval pilot on one GPT-5.4 candidate per task
(155 correct, 9 incorrect) calibrates the setup: self-check accepts $6/9$ incorrect
candidates, structured steelman evaluation accepts $3/9$, and an execution-grounded
control running the reference tests rejects all $9$. On 947 MBPP tasks across five
model families --- 5$\times$947 = 4,735 nominal candidate pairs, of which 4,718
were actually stored (947+938+942+944+947); the 17-pair gap is one-directional and
is not a hidden retraction: it is the zero-stored cells (9 for
ds-v4-flash, 5 for glm-5.3-flash, 3 for qwen38-flash; the two GPT families
store the full 947 each) --- we run self-check and steelman evaluation on
every stored candidate under a complete-response protocol that preserves
raw responses before parsing and keeps malformed records in the denominators,
with a neutral re-examination prompt at matched token use as control. Steelman lowers false acceptance in every family,
surviving Holm adjustment for Qwen ($p=0.0117$) and GPT-6-Luna
($p=2.1\times10^{-6}$); the neutral prompt does not reproduce Qwen's reduction, and
for GPT-6-Luna the two prompts accept incorrect candidates at nearly the same rate.
Each family also evaluates the other families' candidates: acceptance stays elevated
for same-family evaluators after evaluator and generator fixed effects, and the
effect survives a same-provider control. A severity audit of 1,987 incorrect
candidates finds $84.4\%$ fail every reference assert: $62.8\%$ on wrong values,
$37.0\%$ in the crash-or-guard-block bucket, and $0.2\%$ with a missing or renamed
entry point. Consistent with the amplification proposition,
repeated unanimous acceptance on MBPP leaves $27$--$44\%$ residual error, and on
MBPP accepted candidates remain only $55$--$70\%$ correct.
\end{abstract}

\begin{IAENGkeywords}
LLM, self-correction, self-evaluation, information theory, error correlation, external selection, benchmark evaluation.
\end{IAENGkeywords}

\IAENGpeerreviewmaketitle

\section{Introduction}

Large language models cannot reliably self-correct their reasoning~\cite{huang2024selfcorrect}.
Tsui et al. argue that the failure is not a knowledge deficit: their models correct
identical errors presented as external input while failing on the same errors in their
own outputs, a phenomenon termed the \textbf{Self-Correction Blind Spot} and measured
at an average 64.5\% failure rate across 14 models~\cite{tsui2025selfcorrection}.
Our own cross-evaluation finds permissiveness to be largely general, so this study
treats that account as motivation rather than as a premise. This paper argues that the failure is a
structural property of certain evaluation configurations: when generator and evaluator
share failure modes, agreement between them supplies weak evidence of correctness. A
derivation can be elegant, internally consistent, and convincingly presented, yet contain
an error the model cannot detect in its own output.

Three questions are routinely conflated in this setting: whether a model can generate a
correct answer, whether it can recognize its own errors under an internal evaluation
procedure, and how its selection behavior changes when an external correctness channel is
introduced. The first is a generation problem; the second is the subject of this paper;
the third is where the remedy lies. The distinction matters because improvements in
generation do not automatically improve recognition, and because what the evaluator is
asked to do---not simply how much computation it spends---determines what an acceptance
verdict is worth.

\subsection{Contributions}

\begin{enumerate}
    \item \textbf{A measurable diagnostic.} We define the false-acceptance rate $\kappa$,
    true-acceptance rate $\tau$, discrimination gap $\Delta=\tau-\kappa$, and
    acceptance log-likelihood evidence $E=\log_2(\tau/\kappa)$ on frozen candidates, and
    show how to estimate them without any candidate revision, resampling, or repair
    (Section~\ref{sec:formalization}).
    \item \textbf{Paired measurements at two scales.} On the complete 164-task HumanEval
    benchmark, one frozen GPT-5.4 candidate per task yields false acceptance of $6/9$,
    $3/9$, and $0/9$ under self-check, steelman, and an execution-grounded positive
    control. On 947 MBPP tasks across five model families, steelman lowers false
    acceptance in every family, with five-family Holm-adjusted significance for Qwen and
    GPT-6-Luna. Supplementary arms---repeated evaluation, candidate replication, a
    registered moderation analysis, a neutral re-examination control, and a full
    cross-evaluation matrix---localize what the improvement is and is not
    (Section~\ref{sec:empirical}).
    \item \textbf{Conditional theoretical bounds.} Under a shared-blind-spot modeling
    assumption, self-evaluation is bounded in what it can add about correctness
    (Theorem~\ref{thm:info_bound}); $k$ rounds of self-critique inherit the same ceiling
    (Lemma~\ref{lem:repeated}); naive independence treatment of coupled evaluations can
    amplify confidence in errors (Proposition~\ref{prop:miscalibration}); and a selector
    satisfying measurable acceptance-rate bounds supplies a quantifiable evidence floor
    (Proposition~\ref{prop:multi_agent}). The proofs are short; the contribution is the
    measurement program they generate, and the central open problem---an independently
    specified blind-spot variable $Z$---is stated precisely with concrete candidate
    approaches.
    \item \textbf{An auditable protocol and design principles.} Every response is
    preserved before parsing, no output cap is imposed, malformed responses remain
    visible with fixed denominators, and all data artifacts carry SHA-256 manifests in a
    public repository. The design principles for external selection are offered as
    engineering heuristics motivated by the analysis, not as theorem-derived guarantees.
\end{enumerate}

\paragraph{Scope.} This is an \emph{empirical measurement paper with a conditional
theoretical framework}, not an impossibility paper. The information bounds are standard
chain-rule and data-processing consequences whose role is to organize the measurable
quantities, not to constitute a new limitation theorem. The central negative results
assume $S\perp T\mid(G,Z)$: a specific failure regime in which the selector's evaluation
is so corrupted by shared blind spots that the input itself supplies no residual
correctness signal. Selectors with an independent ground-truth channel---a proof checker,
test executor, or independently trained model---fall outside that regime. Cross-generator
error agreement is independently documented at scale~\cite{kim2025correlated}, though
agreement between generators is not itself evidence of evaluator false acceptance. We do
not claim all self-evaluation fails, only that it can fail in this regime, and we measure
when it does.

\section{The Measurement Framework}
\label{sec:formalization}

This section defines the objects every later claim uses. Let the input space be
$\mathcal{X}$ and hypothesis space $\mathcal{H}$. For an input $X\sim\mathcal{D}$:

\begin{Definition}[Generator and Selector]
A \textbf{generator} produces candidate hypotheses via a conditional distribution
$G \sim p(g \mid x)$. A \textbf{selector} produces an evaluation score via
$S \sim p(s \mid x, g)$. The \textbf{acceptance decision} is a deterministic threshold:
\begin{equation}
A := \mathbf{1}\{S \geq \theta_A\}
\end{equation}
\end{Definition}

\begin{Definition}[Correctness Indicator]
Let $C:\mathcal{X}\times\mathcal{H}\to\{0,1\}$ be a task-specific correctness
criterion fixed independently of the selector's acceptance decision. We define the
\textbf{correctness indicator}:
\begin{equation}
T := C(X,G).
\end{equation}
Thus $T=1$ when the generated candidate satisfies the relevant correctness criterion and
$T=0$ otherwise. Exact-answer settings are recovered by taking
$C(x,g)=\mathbf{1}\{g=H^\star(x)\}$ for a target function $H^\star$; executable
benchmarks may instead use a pre-specified behavioral criterion, such as passing a fixed
reference-test suite. The formal results below require only the binary variable $T$, not
uniqueness of a correct hypothesis.
\end{Definition}

\begin{Definition}[Error Events]
The \textbf{generator error event} is $E_G := \{T = 0\}$. The \textbf{selector error
event} is:
\begin{equation}
E_S := \{A \neq T\}
\end{equation}
That is, $E_S$ occurs when the selector accepts an incorrect hypothesis ($A=1, T=0$) or
rejects a correct one ($A=0, T=1$).
\end{Definition}

\begin{Definition}[Conditional Error Coupling]
The \textbf{conditional error coupling} is:
\begin{align}
\kappa &:= \Pr(E_S = 1 \mid E_G = 1) = \Pr(A = 1 \mid T = 0), \\
\tau &:= \Pr(A = 1 \mid T = 1).
\end{align}
Here $\kappa$ is the probability that the selector accepts given that the generator is wrong,
and $\tau$ is the true-acceptance rate. Strong
coupling ($\kappa \approx 1$) indicates that the selector fails to detect generator errors.
Operationally $\kappa$ is a \emph{false-acceptance rate}, not a measured correlation of
failure mechanisms: an evaluator that always accepts has $\kappa=1$ regardless of whether
it shares any latent blind spot with the generator. We retain ``coupling'' as the
framework's term for the conditional quantity, but the experiments estimate evaluator
discrimination, not the latent variable $Z$ or the assumption
$S \perp T \mid (G,Z)$; attributing high $\kappa$ to a shared-blind-spot mechanism is a
hypothesis the framework supplies, not something the measurements identify.
\end{Definition}

Two companion quantities complete the diagnostic. The \textbf{discrimination gap}
$\Delta=\tau-\kappa$ measures how well the evaluator separates correct from incorrect
candidates, and the \textbf{acceptance log-likelihood evidence}
$E=\log_2(\tau/\kappa)$ measures the bits of evidence about correctness supplied by an
acceptance verdict---an event-specific likelihood ratio, not a mutual information.
Overall accuracy is deliberately not the endpoint: on a corpus with 155 correct and 9
incorrect candidates, an always-accept evaluator scores $94.5\%$ while detecting no
errors at all.

\paragraph{The frozen-candidate principle.} All experiments below evaluate one immutable
candidate per task, hashed before any evaluation, with ground truth assigned once by the
benchmark's reference tests and never recomputed. Evaluator prompts forbid modification,
repair, alternative solutions, and execution. The experiment therefore measures evaluator
acceptance, not correction success, and no candidate property can confound a within-task
arm comparison.

\subsection{Shared Blind Spots}

To model one mechanism by which errors might correlate, we introduce a latent variable
capturing shared failure modes.

\begin{Definition}[Blind Spot Variable]
Let $Z \in \mathcal{Z}$ be a latent variable indexing input regions where a model family
systematically struggles. We say generator and selector \textbf{share blind spots} indexed
by $Z$ if there exists a set $\mathcal{Z}_{\mathrm{bad}} \subset \mathcal{Z}$ such that:
\begin{align}
\Pr(T = 0 \mid Z \in \mathcal{Z}_{\mathrm{bad}}) &\geq \alpha_G \gg \Pr(T = 0) \\
\Pr(A = 1 \mid T = 0, Z \in \mathcal{Z}_{\mathrm{bad}}) &\geq \alpha_S > \Pr(A = 1 \mid T = 0)
\end{align}
while both failure rates are low on $\mathcal{Z} \setminus \mathcal{Z}_{\mathrm{bad}}$.
\end{Definition}

Shared blind spots arise naturally when generator and selector share training data,
architecture, or context. For the formal results to have practical force, $Z$ needs an
independent, a priori characterization---a fixed taxonomy of known failure modes for a
model family---not a post hoc fit to observed joint failure. \cite{erroratlas2026} provide one candidate starting point: ErrorAtlas, an empirically
derived taxonomy built from 83 LLMs across 35 datasets. The published atlas can be
frozen before applying it to the generator--selector pairs analyzed here. The Self-Correction Blind Spot of \cite{tsui2025selfcorrection} is consistent with
high $\alpha_S$ in the self-evaluation setting.

\begin{Remark}[The modeling assumption]
\label{rem:ci_assumption}
The selector $S \sim p(s \mid x, g)$ has access to the input $x$ as well as the generator output $g$, and since $T=C(X,G)$ generally depends on $x$, the selector could in principle be informative about $T$ by comparing $g$ against the correct answer for that $x$, a channel independent of $Z$. The assumption $S \perp T \mid (G, Z)$ therefore encodes a specific failure regime: the selector's evaluation process is itself corrupted by the same shared blind spots $Z$, so accessing $x$ yields no correctness signal beyond what $Z$ already mediates. This is precisely the self-evaluation setting, where a model evaluating its own output brings the same failure modes to the same inputs. A proof checker, test executor, or independently trained model violates the assumption only if it supplies residual correctness information, $I(T;S\mid G,Z)>0$. Independent ground-truth access alone does not guarantee this condition: the channel may be unused or uninformative given $(G,Z)$. When residual information is present, external selection can add information beyond the shared blind spots. All subsequent results are conditional on this assumption: an empirically plausible failure regime, not a universal property of self-evaluation.
\end{Remark}

\section{Conditional Bounds on Self-Evaluation}
\label{sec:self_eval_fails}

The central quantity is $I(T;S\mid G)$: the information the selector provides about
correctness given that we already observe the generator output. Conditioning on $G$
rather than on $(G,X)$ is a deliberate modeling choice: under
Remark~\ref{rem:ci_assumption}'s assumption, the selector's access to $X$ yields no
correctness signal beyond what $Z$ mediates, and selectors whose use of $X$ supplies
residual correctness information given $(G,Z)$ fall outside this section's scope. Two
theorems formalize distinct consequences: an information-theoretic ceiling
(Theorem~\ref{thm:info_bound}) and an acceptance evidence bound
(Theorem~\ref{thm:evidence_bound}). Correlated errors can make self-evaluation weakly
informative, but positive association alone does not imply negligible evidence: the
evidence bound is determined by measurable acceptance rates, whereas the ceiling
additionally requires the latent-variable assumptions.

\subsection{Information-Theoretic Scaffolding: Bounding What $Z$ Mediates}

The first theorem bounds what shared blind spots permit any selector to reveal.

\begin{Theorem}[Shared-Blind-Spot Information Bound]
\label{thm:info_bound}
Let $Z$ be a latent variable. Assume the conditional independence
\begin{equation*}
S \perp T \mid (G, Z).
\end{equation*}
Then:
\begin{equation}
I(T; S \mid G) \leq I(T; Z \mid G)
\end{equation}
In particular, if $T \perp Z \mid G$, then $I(T; S \mid G) = 0$.
\end{Theorem}

\begin{proof}
By the chain rule for conditional mutual information:
\begin{equation}
I(T; S, Z \mid G) = I(T; Z \mid G) + I(T; S \mid G, Z)
\end{equation}
Under $S \perp T \mid (G, Z)$, we have $I(T; S \mid G, Z) = 0$, hence:
\begin{equation}
I(T; S, Z \mid G) = I(T; Z \mid G)
\end{equation}
Also by the chain rule:
\begin{multline}
I(T; S, Z \mid G) = I(T; S \mid G) + I(T; Z \mid G, S) \\
\geq I(T; S \mid G)
\end{multline}
since mutual information is nonnegative. Combining yields $I(T; S \mid G) \leq I(T; Z \mid G)$.
If additionally $T \perp Z \mid G$, then $I(T; Z \mid G) = 0$ and the result follows.

\end{proof}

\begin{Remark}[On the informativeness of the bound]
\label{rem:vacuity}
The bound is non-vacuous only under two conditions: (1)~$Z$ is constrained
independently of $T$ and $S$---without this, $Z$ could always be chosen post hoc to encode
all residual dependence between $S$ and $T$, rendering the bound trivially satisfied---and
(2)~$I(T; Z \mid G)$ is shown to be small, so the cap is tight; if $Z$ carries substantial
information about correctness beyond what $G$ reveals, the bound permits $S$ to be highly
informative. Theorem~\ref{thm:info_bound} therefore identifies a structural
relationship---information from $S$ is $Z$-mediated---rather than proving that
self-evaluation is empirically weak. Practical application requires an a priori
characterization of $Z$ (e.g., a fixed taxonomy of known failure modes) that does not
depend on the observed outcome; the paper does not independently establish condition (2)
for real LLM systems. Lemma~\ref{lem:repeated} and Proposition~\ref{prop:multi_agent}
carry more direct practical content because their conditions are expressible in directly
measurable probabilities.
\end{Remark}

\subsection{The Necessity of Falsifiable $Z$}
\label{sec:falsifiable_z}
As Remark~\ref{rem:vacuity} notes, $Z$ can always be constructed post hoc, so
Theorem~\ref{thm:info_bound} cannot by itself predict or diagnose failure in real
systems; it becomes empirically meaningful only when $Z$ is characterized independently
of the failure events being analyzed---defined prior to observation, estimable from
separate data, and constrained to a fixed taxonomy. The theorem's contribution is to
identify exactly what must be measured to determine whether a system is in the failure
regime. The ontological claim that some such $Z$ exists is guaranteed by construction
and carries no empirical content; only the epistemological claim that $Z$ can be
independently specified gives the framework diagnostic power. Candidate routes to an
independently specified $Z$ include behavioral clustering of error types, probing
studies of model internals, and difficulty-stratified benchmarks; the ErrorAtlas
taxonomy shows such a priori specifications are achievable. Constructing a $Z$ that is
both independently specified and informative enough to make
Theorem~\ref{thm:info_bound} non-vacuous remains the key open empirical problem.

The remaining results carry more direct practical content because their conditions are
expressible in measurable probabilities. \cite{lin2024criticbench}'s CriticBench
evaluates the same model on generation, binary critique, and correction for identical
problems, so a pre-specified binary acceptance rule could recode its critique outcomes
into estimates of $\Pr(A{=}1\mid T{=}0)$ and $\Pr(A{=}1\mid T{=}1)$; similarly,
\cite{kadavath2022ptrue}'s P(True) scores could be thresholded at a pre-fixed
$\theta_A$ to construct such rates. Neither source itself defines $\kappa$ or $\tau$.

\begin{Lemma}[Post-Processing Cannot Add Evidence]
\label{lem:post_processing}
For any deterministic acceptance rule $A = \mathsf{acc}(S)$:
\begin{equation}
I(T; A \mid G) \leq I(T; S \mid G)
\end{equation}
\end{Lemma}

This follows directly from the data processing inequality. The acceptance decision cannot
contain more information than the selector score from which it derives.

\subsection{Evidence Bound Formulation}

\begin{Theorem}[Bounded Evidence from Acceptance]
\label{thm:evidence_bound}
Assume the selector has high false acceptance rate:
\begin{equation}
\Pr(A = 1 \mid T = 0) \geq 1 - \varepsilon
\end{equation}
and $\Pr(A = 1 \mid T = 1) \leq 1$. Then the log-likelihood ratio contributed by observing
$A = 1$ satisfies:
\begin{equation}
\log_2 \frac{\Pr(A = 1 \mid T = 1)}{\Pr(A = 1 \mid T = 0)} \leq \log_2 \frac{1}{1 - \varepsilon}
\end{equation}
\end{Theorem}

\begin{proof}
We have $\Pr(A = 1 \mid T = 1) \leq 1$ and $\Pr(A = 1 \mid T = 0) \geq 1 - \varepsilon$,
so:
\begin{equation}
\frac{\Pr(A = 1 \mid T = 1)}{\Pr(A = 1 \mid T = 0)} \leq \frac{1}{1 - \varepsilon}
\end{equation}
Taking $\log_2$ gives the result. 
\end{proof}

\begin{Corollary}[Degenerate Evidence for Small $\varepsilon$]
For small $\varepsilon$:
\begin{equation}
\log_2 \frac{1}{1 - \varepsilon} \approx \frac{\varepsilon}{\ln 2} \approx 1.44\varepsilon
\text{ bits}
\end{equation}
To illustrate the mathematical behavior of the bound: for $\varepsilon = 0.01$ (selector
accepts 99\% of incorrect hypotheses), acceptance provides at most $0.014$ bits of evidence.
For $\varepsilon = 0.35$ (a more moderate regime), the exact bound is $0.62$ bits
(the small-$\varepsilon$ approximation gives $0.50$). These are mathematical illustrations of how the bound scales with $\varepsilon$, not
empirical estimates; determining which regime applies to a given deployed system requires
direct measurement of acceptance probabilities.
\end{Corollary}

Theorem~\ref{thm:evidence_bound} provides an \emph{upper} bound on evidence under high
false acceptance; outside that regime acceptance may convey substantial evidence, and
many successful LLM applications violate the preconditions by incorporating external
feedback channels. The regime is directly measured: the MBPP point estimates
$\hat\kappa=0.86$--$0.97$ of Section~\ref{sec:empirical} fall inside it; the HumanEval
estimates do not, and the bound is not applied to them.

\subsection{Sufficient Conditions for Positive Evidence}

\begin{Proposition}[Conditions for positive evidence]
\label{prop:multi_agent}
Suppose a selector satisfies
\begin{equation*}
\Pr(A = 1 \mid T = 0) \leq 1 - \epsilon,\qquad
\Pr(A = 1 \mid T = 1) \geq \delta
\end{equation*}
for some $\delta > 1 - \epsilon > 0$. Then acceptance by that selector provides
\begin{equation}
\log_2 \frac{\delta}{1 - \epsilon} > 0
\end{equation}
bits of evidence about correctness.
\end{Proposition}

\begin{proof}
The log-likelihood ratio of $A=1$ is
$\log_2 \frac{\Pr(A=1 \mid T=1)}{\Pr(A=1 \mid T=0)} \geq \log_2 \frac{\delta}{1-\epsilon}$,
which is positive whenever $\delta > 1 - \epsilon$. 
\end{proof}

\begin{Remark}
The proposition gives \emph{sufficient} conditions under which acceptance provides positive
evidence, not necessary and sufficient conditions. Positive evidence requires only
$\Pr(A=1 \mid T=1) > \Pr(A=1 \mid T=0)$; the $(\delta, \epsilon)$ parameterization makes
the bound explicit and measurable. A selector meeting these conditions provides at least
$\log_2(\delta/(1-\epsilon))$ bits; the actual evidence may be larger. The proposition
additionally requires the selector to accept correct hypotheses at rate $\geq \delta >
1-\epsilon$, since a trivially rejecting selector would satisfy the false-acceptance bound
without being useful.
\end{Remark}

\subsection{Average Information and Naive-Confidence Miscalibration}

The average-information ceiling and confidence miscalibration are different phenomena;
the second requires an explicit observer model.

\begin{Lemma}[Repeated Self-Critique Bound]
\label{lem:repeated}
Consider $k$ selector outputs $S_1, \ldots, S_k$ with acceptance decisions
$A_i = \mathbf{1}\{S_i \geq \theta_A\}$. Assume $(S_1, \ldots, S_k) \perp T \mid (G, Z)$
jointly. A sufficient condition is that the $S_i$ are mutually conditionally independent
given $(T,G,Z)$, with each conditional law $\Pr(S_i\mid T,G,Z)$ independent of $T$.
The resulting product is then free of $T$ for every $(G,Z)$, which implies the required
joint conditional independence. Conditional independence given only $(G,Z)$, together with
each separate condition $S_i\perp T\mid(G,Z)$, would not suffice. Then:
\begin{equation}
I(T; A_{1:k} \mid G) \leq I(T; Z \mid G)
\end{equation}
That is, $k$ critiques provide no more information about correctness than the single blind
spot variable $Z$.
\end{Lemma}

\begin{proof}
By the chain rule for conditional mutual information:
\begin{equation}
I(T; A_{1:k}, Z \mid G) = I(T; Z \mid G) + I(T; A_{1:k} \mid G, Z)
\end{equation}
Under the joint conditional independence assumption, $A_{1:k} \perp T \mid (G, Z)$, so
$I(T; A_{1:k} \mid G, Z) = 0$, giving $I(T; A_{1:k}, Z \mid G) = I(T; Z \mid G)$.
Also by the chain rule:
\begin{multline}
I(T; A_{1:k}, Z \mid G) = I(T; A_{1:k} \mid G) \\
+ I(T; Z \mid G, A_{1:k}) \geq I(T; A_{1:k} \mid G)
\end{multline}
since mutual information is nonnegative. Combining yields $I(T; A_{1:k} \mid G) \leq I(T; Z \mid G)$. 
\end{proof}

\begin{Remark}
The joint independence assumption is strictly stronger than the individual conditions
$S_i \perp T \mid (G,Z)$; individual conditional independence
does not in general imply the joint statement. Fresh context alone does not establish the
lemma's factorization sufficient condition.
\end{Remark}

\begin{Remark}[The ceiling bounds total information]
\label{rem:marginal}
Lemma~\ref{lem:repeated} should not be read as ``repeated evaluation is useless.'' The
chain rule decomposes the bounded quantity into marginal contributions:
\begin{multline*}
I(T; A_{1:k} \mid G) \\
= \sum_{i=1}^{k} I(T; A_i \mid G, A_{1:i-1}),
\end{multline*}
and the lemma constrains the sum, not the individual terms. Each critique can contribute
substantial positive information on the way toward the ceiling; only the cumulative total
is capped at $I(T; Z \mid G)$. Concretely, let $T = Z$ be fair binary with $G$
uninformative about $T$, so the ceiling is $I(T;Z\mid G) = 1$ bit, and let each critique
be an independent binary-symmetric reading of $Z$ with accuracy $0.7$. The conditional
independence assumption holds, a single critique carries $\approx 1 - h_2(0.3) \approx 0.12$
bits, subsequent marginal gains shrink as the ceiling is approached, and ten critiques
carry $I(T;A_{1:10}) \approx 0.678$ bits (computed exactly from the binomial mixture),
not the full bit. The lemma remains true throughout: the total
never exceeds the ceiling.
\end{Remark}

\begin{Remark}[Average versus event-specific evidence]
\label{rem:not_llr}
The bounded quantity in Lemma~\ref{lem:repeated} is an \emph{average} mutual
information. It is not a bound on the log-likelihood ratio of any particular sequence
of decisions, and it does not follow that a confident posterior after repeated
acceptance is unjustified. In the same binary-symmetric example, ten unanimous YES
verdicts have likelihood ratio $(0.7/0.3)^{10}$, i.e.\ a log-likelihood ratio of
$\approx 12.22$ bits, while $I(T;A_{1:10}) \approx 0.678$ bits. Both statements are
true simultaneously: the posterior update on that particular sequence is large and
statistically legitimate, because the readings really are conditionally independent
given $T$ (here $T=Z$). What the lemma bounds is the information averaged over all
$2^{10}$ verdict patterns, most of which are not unanimous. Consequently,
Lemma~\ref{lem:repeated} on its own does \emph{not} establish false confidence
amplification. Three distinct claims must be separated: (i) bounded average mutual
information, which the lemma proves; (ii) dependence between evaluation outcomes,
which is a property of the sampling model; and (iii) miscalibration caused by
treating dependent evaluations as independent, which requires an explicit
miscalibration statement and is the content of
Proposition~\ref{prop:miscalibration} below.
\end{Remark}

\begin{Proposition}[Naive-independence miscalibration]
\label{prop:miscalibration}
This is a separate existence result, not a consequence of Lemma~\ref{lem:repeated}.
It assumes evaluations are i.i.d.\ given $(T,Z)$, whereas the lemma assumes that the
evaluation vector is conditionally independent of $T$ given $(G,Z)$. These are distinct
sampling assumptions. Fix $G$ and suppress it, and define
$a_z = \Pr(A_i{=}1 \mid T{=}1, Z{=}z)$ and $b_z = \Pr(A_i{=}1 \mid T{=}0, Z{=}z)$,
and let $\rho = \Pr(T{=}1) \in (0,1)$. After $k$ unanimous acceptances, the true
posterior odds are
\[
O_{\mathrm{true}}(k) \;=\; \frac{\rho}{1-\rho}\,
\frac{\mathbb{E}[a_Z^{\,k} \mid T{=}1]}{\mathbb{E}[b_Z^{\,k} \mid T{=}0]},
\]
whereas an observer who wrongly treats the evaluations as independent given $T$
alone, using the marginal rates $\bar a = \mathbb{E}[a_Z \mid T{=}1]$ and
$\bar b = \mathbb{E}[b_Z \mid T{=}0]$, computes
\[
O_{\mathrm{naive}}(k) \;=\; \frac{\rho}{1-\rho}\left(\frac{\bar a}{\bar b}\right)^{k}.
\]
The two posteriors can disagree, but non-degeneracy alone does not imply disagreement:
if $a_Z$ and $b_Z$ have identical distributions under their respective correctness
classes, the odds coincide for every $k$ even when both are non-degenerate. Exponential divergence requires additional
conditions. In particular, take $\bar a > \bar b > 0$ with
$a_z \equiv a < 1$ constant and, under $T{=}0$, $b_Z = 1$ with probability $q$ (a
shared blind spot: every evaluation accepts) and $b_Z = 0$ otherwise, so
$\bar b = q < a$. Then
\begin{gather*}
\log O_{\mathrm{naive}}(k) = \log\tfrac{\rho}{1-\rho} + k\log\tfrac{a}{q}
\;\longrightarrow\; +\infty, \\
\log O_{\mathrm{true}}(k) = \log\tfrac{\rho}{1-\rho} + k\log a - \log q
\;\longrightarrow\; -\infty .
\end{gather*}
The naive observer becomes arbitrarily confident that the output is correct while
the true posterior odds of correctness vanish.
\end{Proposition}

\begin{proof}
Conditional on $(T,Z)$ the $k$ acceptances are independent, so the likelihood of the
unanimous sequence is $a_Z^k$ under $T{=}1$ and $b_Z^k$ under $T{=}0$; taking
expectations over $Z$ gives $O_{\mathrm{true}}$. The naive observer replaces each
$k$-th moment by the $k$-th power of a mean. In the two-point example,
$\mathbb{E}[a_Z^k \mid T{=}1] = a^k$ and
$\mathbb{E}[b_Z^k \mid T{=}0] = q\cdot 1^k + (1-q)\cdot 0 = q$ for every $k$, so
$O_{\mathrm{true}}(k) = \frac{\rho}{1-\rho} a^k / q$; the limits follow from
$a/q > 1$ and $0 < a < 1$.
\end{proof}

\paragraph{Intuition: what the failure mode actually is.}
The failure is not that confident posteriors after repeated acceptance are inherently
unjustified---Remark~\ref{rem:not_llr} gives a counterexample in which they are
justified. The failure is a specific accounting error: an observer who treats
$Z$-coupled evaluations as independent confirmations substitutes powers of means for
means of powers, and in the blind-spot regime this substitution does not merely
overshoot but flips the sign of the evidence. The subjective experience is $k$
confirmations; the true likelihood of that unanimous sequence is dominated by a single
shared cause whose probability does not shrink with $k$. Whether real evaluator
populations sit in this sign-reversal regime is an empirical question; the experiments
below measure the discrimination quantities $\kappa$ and $\tau$ that enter it, not $Z$
itself.

\paragraph{What the theory buys the experiments.} The proofs in this section are short;
their function is to organize the measurement program of Section~\ref{sec:empirical}.
The repeated-critique ceiling frames the repeated-evaluation and replication arms; the
conditional-independence assumption marks what the moderation analysis and
cross-evaluation matrix can and cannot reach; the evidence bound fixes the reporting of
$E$ alongside the rates; and the necessity of falsifiable $Z$ marks the boundary beyond
which these experiments cannot see. The theory organizes the measurements and identifies
their limits; it is not what they test.

\section{Experiments}
\label{sec:empirical}

\subsection{Protocol and Reproducibility}
\label{sec:protocol}

Both experiments instantiate the frozen-candidate principle of
Section~\ref{sec:formalization}. Ground truth $T$ is assigned once by the benchmark's
reference tests before any evaluation; every candidate is deterministically hashed and
frozen; no candidate is selected, repaired, revised, or resampled as a function of
correctness; and evaluator prompts forbid modification, repair, alternative solutions,
and execution. The three evaluator conditions are: \textbf{self-check} (the same model
judges its frozen candidate in a fresh context with a binary verdict),
\textbf{steelman} (the same model produces the strongest case against and for the
candidate before a final verdict), and an \textbf{execution-grounded positive control}
(the evaluator receives only a binary PASS/FAIL signal from the reference-test
mechanism that defines $T$---an oracle-grounded control, not an independent validation
benchmark).

The MBPP runs use a complete-response protocol: every API attempt is appended verbatim,
before any parsing, to an immutable JSONL event ledger recording model, stage, task,
request parameters, raw response text, provider metadata, and \texttt{finish\_reason};
final attempts carry no \texttt{max\_tokens}-style parameter, and no response entering
any denominator is truncated at our behest; a pair whose final attempt reaches the provider's own output ceiling is
never regenerated and stays visible as a malformed record with a fixed denominator;
verdict extraction accepts only a strict terminal \texttt{VERDICT: ACCEPT} or
\texttt{VERDICT: REJECT} marker with no substring fallback. Across the five ledgers
(production, repeated evaluation, replication, neutral control, cross-evaluation),
93,261 API events are preserved, and every parsed verdict reproduces from preserved raw
text. All generations and evaluations ran at temperature 0; GPT-5.4 used the same
pinned snapshot as the HumanEval experiment (gpt-5.4-2026-03-05), and model identifiers
and request parameters are recorded per event. All 164 archived HumanEval steelman
prompts contain the no-repair and no-corrected-code instructions, and inspection of all
code-fenced responses found candidate quotations or illustrative expressions, not a
replacement implementation.

The preserved harness, candidates, raw responses, parsed decisions, validators, and
deterministic analysis scripts are public at
\url{https://github.com/AndBrilliant/selfcorr-benchmarks} (commit
\texttt{81537cf}, the frozen data snapshot); the audited release adds SHA-256
manifests covering every released data artifact. The HumanEval contrasts
reproduce from \path{scripts/humaneval/paired_contrasts.py}, the MBPP analyses from
\path{scripts/mbpp_v2_analysis.py} and its companion scripts, and ledger integrity from
\path{scripts/validate_mbpp_v2_pilot.py}.

\subsection{HumanEval: Measuring the Diagnostic}

The complete 164-task HumanEval benchmark~\cite{chen2021codex} supplies executable
reference tests and a deterministic functional-correctness criterion. GPT-5.4
(``gpt-5.4-2026-03-05'', temperature 0, reasoning effort ``none'') produced one frozen
candidate per task: 155 correct, 9 incorrect.\footnote{Baseline generation was
interrupted after initiating the request for HumanEval/85 but before any response was
persisted. On resume, that task remained unfrozen and the identical request was
reissued. No candidate content from the interrupted request was observed or used in the
decision to reissue it.} The same 164 candidates were judged under
self-check, steelman, and the execution-grounded control, all in fresh contexts. Wilson
score intervals are used for the binomial rates, and marginal uncertainty in $\hat E$ is
obtained by transforming a two-sided 95\% Koopman score interval for $\tau/\kappa$ with
the Miettinen--Nurminen correction~\cite{koopman1984ratio,miettinen1985comparative};
no pseudocount is added.

\begin{table*}[tp]
\caption{Paired evaluator acceptance on the identical frozen HumanEval corpus
($N=164$: 155 correct, 9 incorrect). Wilson 95\% confidence intervals are shown for
$\kappa$ and $\tau$. Log-likelihood evidence from acceptance is $E=\log_2(\tau/\kappa)$;
its marginal intervals are transformed two-sided 95\% Koopman score intervals with the
Miettinen--Nurminen correction. The execution arm has zero observed false acceptances, so
its plug-in $\hat E$ and the upper interval endpoint are unbounded.}
\label{tab:results}

\begin{tabular}{l c c c c c}
\hline
\textbf{Evaluator} & \textbf{Wrong accepted} & $\hat{\kappa}$ [95\% CI] &
\textbf{Correct accepted} & $\hat{\tau}$ [95\% CI] & $\hat{E}$ (bits) [95\% CI] \\
\hline
Self-check & 6/9 & 0.667 [0.354, 0.879] & 146/155 & 0.942 [0.893, 0.969] & 0.50 [0.09, 1.42] \\
Steelman & 3/9 & 0.333 [0.121, 0.646] & 149/155 & 0.961 [0.918, 0.982] & 1.53 [0.57, 3.00] \\
Execution-grounded & \textbf{0/9} & \textbf{0.000 [0.000, 0.299]} & 149/155 & \textbf{0.961 [0.918, 0.982]} & \textbf{$\infty$ $[1.68,\infty)$} \\
\hline
\end{tabular}
\end{table*}

The two internal arms are ordered in false acceptance: $\hat\kappa$ falls from $6/9$
to $3/9$ while true acceptance rises from $146/155$ to $149/155$; the execution control
rejects all nine incorrect candidates (zero observed false acceptances). The discrimination gaps are
$\hat\Delta=0.275$, $0.628$, and $0.961$. The oracle arm nevertheless rejects six
PASS candidates, so its verdicts are not a mechanical relabeling of $T$. All six
responses are a bare ``NO'' under a verdict-only prompt, so no rationale is preserved;
five of the six candidates were also rejected by self-check and four by steelman.
The arm thus sometimes overrides an external PASS with the same negative judgment
issued by the internal arms. The steelman
result is not a global shift toward skepticism: self-check and steelman each issued
exactly 152 YES and 12 NO verdicts, but steelman placed its rejections more accurately,
reducing false acceptance from $6/9$ to $3/9$ while reducing false rejection from
$9/155$ to $6/155$. Within the nine incorrect candidates, five are accepted only by
self-check and two only by steelman (exact McNemar $p=0.453$); execution grounding
removed the remaining three steelman false acceptances, and across the endpoints it
corrected 10 self-check evaluator errors while introducing 1 ($p=0.0117$). Evidence from acceptance is
$0.50$ bits under self-check ($[0.09,1.42]$) and $1.53$ bits under steelman
($[0.57,3.00]$); with zero false acceptances in the execution arm, the plug-in ratio is
unbounded and the transformed score interval is $[1.68,\infty)$ bits---a finite lower
endpoint, not a claim of infinite population evidence.

A paired bootstrap (4,000-resample inverse-empirical-CDF percentile resampling within
the nine-error and 155-correct strata, seed 20261104, no
pseudocounts~\cite{efron1979bootstrap}; 100 of the 4,000 resamples yield infinite
$\hat E$, so the finite upper endpoint depends on the quantile convention) gives
steelman-minus-self-check contrasts of
$\Delta\hat\kappa=-0.333$ $[-0.889,0.222]$, $\Delta\hat\tau=+0.019$
$[-0.013,0.052]$, and $\Delta\hat E=+1.03$ bits $[-0.71,3.26]$: the nominal,
unadjusted interval includes no evidence gain, so the nine-error stratum cannot support
a precise between-arm evidence estimate. HumanEval is therefore a clean paired
demonstration of the diagnostic ordering, not a powered effect estimate;
Section~\ref{sec:mbpp} supplies the powered counterpart.

\paragraph{Context-separation comparison.} Because both internal arms use fresh
contexts, they compare evaluator instructions, not context inheritance. Same-context
evaluations replayed the original generation conversation and appended the evaluator
instruction; the steelman variant differs in response-format instructions and parsing,
so the comparison does not isolate context exactly. False acceptance was $6/9$ for
same-context self-check and $4/9$ for same-context steelman, against $6/9$ and $3/9$
fresh; true acceptance was $148/155$ and $148/155$ same-context versus $146/155$ and
$149/155$ fresh. The steelman advantage appears in both contexts under these protocol
variants, and the comparison provides no compelling support for a large fresh-context
benefit while remaining inconclusive about moderate or general effects at nine errors.
Context separation is retained as an engineering heuristic, not a validated mechanism.

\subsection{MBPP: Five-Family Paired Measurement}
\label{sec:mbpp}

The HumanEval experiment pins precision by fixing the candidate corpus; the cost is a
nine-task error stratum. As a complement, we measured the same two internal arms on 947
MBPP tasks~\cite{austin2021mbpp} across five model families: DeepSeek-V4.1-Flash,
GLM-5.3-Flash, Qwen-3.8-Flash, GPT-5.4, and GPT-6-Luna. One greedy candidate per
(task, model) cell was generated, hashed, and frozen; ground truth was assigned before
evaluation; and the frozen candidate was judged in fresh contexts under both arms. The
MBPP steelman prompt is a critique-only variant, not verbatim identical to the
HumanEval prompt; neither evaluator output is executed or used to recompute $T$.
Generation accuracy among candidates with assigned ground truth is 54.3\% (GPT-5.4),
57.8\% (DeepSeek), 67.3\% (GLM), 55.9\% (Qwen), and 54.2\% (GPT-6-Luna), giving
incorrect strata of 308--434 per family. Both arms judge the same immutable candidate,
so the arm comparison is paired at the task level and carries an exact McNemar test;
because each family is judged on its own error distribution, cross-family rate
comparisons remain descriptive, not a capability ranking.

\begin{table*}[!t]
\centering
\caption{Fresh complete-response MBPP paired results, five model families. One greedy
candidate per task, frozen and hashed before evaluation; both arms judge the same
candidate in fresh contexts. ``Paired'' is the number of candidates with assigned
ground truth and a terminal verdict from both arms; $n_{T=0}$ is the incorrect stratum
of that paired set, and $\kappa=P(A^{\mathrm H}{=}1\mid T{=}0)$,
$\tau=P(A^{\mathrm H}{=}1\mid T{=}1)$ are strict-verdict false- and true-acceptance
rates with 95\% Wilson score intervals. $p_{\mathrm{Holm}}$ is the five-family
Holm-adjusted value for the paired incorrect-stratum McNemar contrast. Malformed arms are excluded from paired denominators
and reported as counts in the text, never imputed.}
\label{tab:kappa_mbpp}
\resizebox{\textwidth}{!}{%
\begin{tabular}{lrrrrrr}
\hline
Model & Tasks & Paired & $n_{T=0}$ & $\hat\kappa$ sc/st (\%) & $\hat\tau$ sc/st (\%) & $p_{\mathrm{Holm}}$ \\
\hline
DeepSeek-V4.1-Flash (cloud) & 947 & 917 & 389 & 96.9 [94.7,98.2] / 94.3 [91.6,96.2] & 98.9 [97.8,99.4] / 98.9 [97.8,99.4] & 0.261 \\
GLM-5.3-Flash (cloud) & 947 & 917 & 299 & 89.0 [84.9,92.0] / 85.6 [81.2,89.1] & 97.6 [96.0,98.5] / 97.6 [96.0,98.5] & 0.328 \\
Qwen-3.8-Flash (cloud) & 947 & 940 & 413 & 94.9 [92.4,96.7] / 90.6 [87.4,93.0] & 99.2 [98.1,99.7] / 97.3 [95.4,98.3] & $1.2\times10^{-2}$ \\
GPT-5.4 (cloud) & 947 & 938 & 427 & 90.2 [87.0,92.6] / 88.5 [85.2,91.2] & 97.1 [95.2,98.2] / 95.1 [92.9,96.7] & 0.392 \\
GPT-6-Luna (cloud) & 947 & 943 & 432 & 93.5 [90.8,95.5] / 86.8 [83.3,89.7] & 96.5 [94.5,97.8] / 93.3 [90.8,95.2] & $2.1\times10^{-6}$ \\
\hline
\end{tabular}}
\end{table*}

Table~\ref{tab:kappa_mbpp} gives the paired results at two orders of magnitude more
error events than the HumanEval diagnostic. In every family the steelman arm accepts
fewer incorrect candidates. The paired contrast is significant for Qwen ($b/c=26/8$,
exact $p=2.9\times10^{-3}$) and GPT-6-Luna ($b/c=32/3$, exact $p=4.2\times10^{-7}$),
and both survive Holm adjustment across the five families ($p_{\mathrm{Holm}}=0.0117$
and $2.1\times10^{-6}$); the other three families show the same direction without
significance. No overall-verdict McNemar contrast is significant in any family (all
$p>0.11$), and no arm-by-correctness interaction reaches significance either
(per-family Fisher homogeneity $p=0.41$--$0.49$), though that test is unpowered for
establishing the absence of an interaction. The paired discrimination-gap change has a
95\% bootstrap interval of $[-4.2,+3.5]$ points (GPT-5.4), $[-0.3,+5.4]$ (DeepSeek),
$[-1.1,+7.9]$ (GLM), $[-0.6,+5.5]$ (Qwen), and $[+0.4,+7.0]$ (GPT-6-Luna), so the
discrimination improvement is established for GPT-6-Luna only; these intervals are
per-family and unadjusted. True-acceptance cost is small and family-dependent: at most
3.2 points (GPT-6-Luna $96.5\%\rightarrow93.3\%$), and zero for DeepSeek and GLM.

The practical reliability of accepted candidates moves correspondingly little: among
candidates an arm accepts, the fraction correct is $55$--$70\%$ and improves by at most
$1.0$ points under steelman (Qwen $57.2\%\rightarrow57.8\%$, GPT-6-Luna
$55.0\%\rightarrow56.0\%$, GPT-5.4 flat at $56.3\%\rightarrow56.2\%$). A statistically
detectable reduction in false acceptance therefore does not translate into a large
improvement in the reliability of accepted solutions: better internal critique is not
reliable validation. Acceptance itself buys almost nothing over the unfiltered pool:
accepted candidates are only $0.8$--$3.2$ points more likely to be correct than the
family's base generation accuracy (correct$/$all 947 tasks). Absolute rates differ from the HumanEval measurement (86--97\%
versus 33--67\% false acceptance); benchmarks, candidate-generating models, and
incorrect-stratum sizes all differ, so only the ordering transfers descriptively, not
the level. In a full-population audit of all 1,987 incorrect candidates, 84.4\% failed
\emph{every} reference assert (\path{scripts/mbpp_v2_error_severity_audit.py}). The
$86$--$97\%$ levels may partly reflect conventions hidden from the evaluator rather
than evaluator blind spots alone.

On the paired incorrect stratum the two arms' joint acceptance is $80.3$--$92.0\%$
across families, but the level alone is not evidence of association: both arms accept
most incorrect candidates, so independence would already predict joint acceptance of
$76.2$--$91.4\%$ from the marginals. The observed excess over the independence baseline
is $+0.6$ to $+4.9$ percentage points per family (DeepSeek $+0.6$, Qwen $+2.7$,
GPT-5.4 $+3.8$, GLM $+4.1$, GPT-6-Luna $+4.9$); only this excess, not the absolute
level, measures association between the arms' decisions, and it identifies neither a
latent $Z$ nor the conditional-independence assumptions of the theory.

\subsection{Repeated Evaluation, Replication, and Moderation}

\paragraph{Repeated evaluation.} Each frozen candidate was judged in three repeated
calls per arm at temperature zero: the production verdict plus two fresh calls,
adding 18,884 events. These are repeated provider calls, not independent draws
from a controlled sampling model. Verdicts are highly
stable: $89.7$--$96.9\%$ of candidates receive a unanimous verdict across three rounds,
and agreement itself is informative: candidates accepted in all three rounds carry
$27$--$44\%$ error, versus $48$--$81\%$ for candidates accepted in only one or two
rounds (Table~\ref{tab:repeated}). Repeated confident acceptance therefore does add
evidence about correctness on a fixed candidate, but the residual error of unanimous
acceptance stays at $27$--$44\%$ at $k{=}3$; for GPT-5.4, three unanimous acceptances
move the error rate from $45.7\%$ at baseline to $43.4\%$. This is an empirical stability property,
not a measured information-theoretic ceiling: the latent bound is not identified, and
these data do not show that further evaluations could not improve discrimination.
Repeated-evaluation denominators are arm-specific: $n$ counts candidates with a parsed
verdict in all three rounds of that arm and does not require a parsed verdict in the
other arm, so it differs from the production paired stratum.

\begin{table}[!t]
\centering
\caption{Repeated evaluation ($k{=}3$ rounds per arm, temperature zero).
$n$ is the number of candidates with a parsed verdict in all three rounds;
$e_{3/3}$ is the error rate among candidates accepted in all three rounds
(with count); $e_{1:2}$ is the error rate among candidates accepted in one
or two rounds (with count).}
\label{tab:repeated}
\begin{tabular}{llrrr}
\hline
Family & Arm & $n$ & $e_{3/3}$ (\%) & $e_{1:2}$ (\%) \\
\hline
GPT-5.4 & self-check & 947 & 43.4 (874) & 72.4 (29) \\
 & steelman & 936 & 43.6 (808) & 48.1 (79) \\
DeepSeek-V4.1-Flash & self-check & 910 & 40.3 (865) & 78.0 (41) \\
 & steelman & 900 & 41.0 (847) & 71.1 (45) \\
GLM-5.3-Flash & self-check & 890 & 27.3 (788) & 65.2 (92) \\
 & steelman & 903 & 26.6 (793) & 67.0 (88) \\
Qwen-3.8-Flash & self-check & 933 & 42.0 (879) & 80.9 (47) \\
 & steelman & 928 & 41.2 (839) & 67.6 (74) \\
GPT-6-Luna & self-check & 941 & 44.1 (866) & 68.1 (47) \\
 & steelman & 946 & 43.0 (821) & 67.1 (70) \\
\hline
\end{tabular}
\end{table}

\paragraph{Candidate replication.} A full replication generated a second greedy
candidate for every (task, model) cell, so ground truth varies within cells. The two
programs are token-identical in $33.9\%$ of cells (family range $2$--$82\%$), and
correctness flips between replicas in $6.2\%$ of cells, so the replication introduces
genuine candidate variation alongside a large identical-replica component whose
agreement is guaranteed by construction. Evaluator verdicts repeat across the two
candidates $89$--$97\%$ of the time, and an arm's error on one candidate predicts that
the other is incorrect with very small Fisher $p$-values. That association must be
benchmarked before it carries mechanistic weight: an always-accept evaluator would show
the identical association whenever incorrectness persists across two attempts at the
same task. Measured against this baseline, the headline association is largely task
difficulty: the baseline $P(T_B{=}0\mid T_A{=}0)$ is already $75$--$99\%$ per family,
and the observed $P(T_B{=}0\mid\text{arm errs on }A)$ sits at or below it in every
cell. The verdict does contribute a small incremental signal beyond the first
candidate's ground truth: conditional on $T_A{=}0$, a candidate the arm correctly
rejects is followed by an incorrect replica only $53$--$93\%$ of the time, versus
$76$--$99\%$ for a wrongly accepted one; the contrast is significant in 3 of the 10
arm-family cells. Joint acceptance of both greedy draws carries the same
error rate as accepting either ($25$--$44\%$ both versus $30$--$45\%$ either). The
replication is therefore evidence of strong task-level persistence of incorrectness
plus a modest evaluator-specific component, not evidence of shared blind spots.

\paragraph{Registered moderation analysis.} We tested whether the arm association
survives conditioning on visible, correctness-blind features. Five moderators (task
provenance, prompt length, candidate length, entry-point arity, concept class) and an
explicit decision rule were registered after production data collection began and
before any stratified statistic was computed
(the specification file \path{scripts/mbpp_v2_moderator_spec.md} in the pinned
repository commit is byte-identical to the registered version). Under the registered
rule the association is persistent in all
five families: the stratified $\phi$ is within $0.03$ of the crude value for
every moderator (crude range $0.15$--$0.59$), and joint acceptance stays at or above
$68\%$ in every stratum with at least ten paired incorrect candidates
(smallest qualifying stratum $n=13$ at $77\%$; the lone exception is a
one-observation stratum at $0\%$). Here $\phi$ is the self-check--steelman verdict association
on the incorrect stratum, not a ground-truth effect. One moderator partially attenuates it: in the text-processing
stratum $\phi$ drops to $0.07$--$0.30$ versus $0.42$--$0.46$ in the remaining
candidates for three families. This rules out the five registered covariate
explanations, not every candidate-difficulty account, and it is associational rather
than a conditional-independence test.

\subsection{Mechanism: Adversarial Content and Self-Leniency}
\label{sec:mechanism}

Two supplementary arms, run under the same complete-response protocol on the same
frozen candidates, test the two natural explanations of the steelman improvement: added
deliberation and self-specificity.

\paragraph{Neutral re-examination control.} One neutral, non-adversarial step-by-step
verification call per frozen candidate---same ledger, same strict verdict grammar, no
output cap, realized completion tokens recorded---separates adversarial content from
added deliberation. Table~\ref{tab:control} shows the result. The control
false-acceptance rate lies between the two production arms for DeepSeek, GLM, and
Qwen, sits at the self-check level for GPT-5.4, and reaches the steelman level for
GPT-6-Luna. The paired control-versus-steelman contrast is significant for exactly one
family, Qwen ($p=0.0052$; Holm-adjusted across the five control contrasts, $p=0.026$;
control using $0.99\times$ the steelman tokens): for Qwen, matched deliberation
without adversarial framing accepts incorrect candidates at the self-check level:
this neutral prompt at matched token use does not reproduce the steelman reduction,
localizing Qwen's gain to the steelman protocol rather than deliberation alone. GPT-5.4
shows the same direction with the computation imbalance reversed---control uses
$1.44\times$ the steelman tokens yet accepts $90.6\%$ of incorrect candidates versus
$88.5\%$ under steelman---but the paired contrast is not significant there ($p=0.14$). For GPT-6-Luna, control and
steelman accept incorrect candidates at essentially the same rate ($85.8\%$ versus
$86.7\%$, both far below self-check $93.5\%$), consistent with content-neutral
re-engagement; that near-equality is not an equivalence result, and a modest
adversarial-content contribution cannot be excluded. Across families the control
falsifies the simplest ``more tokens did it'' account: equal or greater deliberation
reproduces the self-check blind spot in the family with the clearest steelman gain. The
control has measured computation, not fully matched computation (token ratios
$0.65$--$1.44$), and even equal average token use does not establish identical
computation per candidate.

\begin{table}[!t]
\centering
\small
\caption{Neutral re-examination control versus the two production arms on the
triple-paired incorrect stratum (ground truth assigned, parsed verdict from all three
arms; $n_{T=0}$ is its size per family). $\hat\kappa$ is false acceptance in percent;
the token ratio is realized mean control-versus-steelman completion tokens; $b/c$
counts discordant incorrect candidates (control accepts, steelman rejects versus the
reverse); $p$ is the exact McNemar value for the control-versus-steelman contrast.}
\label{tab:control}
\resizebox{\columnwidth}{!}{%
\begin{tabular}{lccccccc}
\hline
Family & $n_{T=0}$ & $\hat\kappa$ self & $\hat\kappa$ steel & $\hat\kappa$ ctrl & tok ratio & $b/c$ & $p$ \\
\hline
GPT-5.4 & 427 & 90.2 & 88.5 & 90.6 & 1.44$\times$ & 19/10 & 0.136 \\
DeepSeek-V4.1-Flash & 386 & 96.9 & 94.3 & 95.1 & 0.65$\times$ & 16/13 & 0.711 \\
GLM-5.3-Flash & 288 & 89.6 & 86.8 & 88.5 & 0.88$\times$ & 21/16 & 0.511 \\
Qwen-3.8-Flash & 410 & 94.9 & 90.7 & 94.6 & 0.99$\times$ & 23/7 & \textbf{0.0052} \\
GPT-6-Luna & 430 & 93.5 & 86.7 & 85.8 & 0.80$\times$ & 12/16 & 0.572 \\
\hline
\end{tabular}}
\end{table}

\paragraph{Cross-evaluation matrix.} Every frozen candidate was additionally judged by
the four \emph{other} families in both arms (18,872 records; 399 malformed, reported
with fixed denominators). Two findings follow. First, permissiveness is largely
general, not self-specific: every evaluator accepts other families' incorrect
candidates at high rates (cross false acceptance $73$--$94\%$ under self-check and
$69$--$92\%$ under steelman), so false acceptance of incorrect code does not require
the code to be one's own. Second, same-family evaluators are more permissive toward their own errors.
Table~\ref{tab:cross_matrix} gives the full matrix. Candidate-paired discordance runs
$2.5$--$7.4$ to one toward external evaluators rejecting errors that the generating
family accepted, versus the reverse, and every task-clustered
paired-bootstrap interval is above zero. A crossed logistic model with categorical
evaluator and generator effects and task-clustered sandwich standard errors separates
this same-family effect from each evaluator's general permissiveness: same-family
acceptance of incorrect candidates remains elevated under self-check
($\beta=1.078$, odds ratio $2.94$, $p=1.2\times10^{-36}$) and steelman
($\beta=0.654$, odds ratio $1.92$, $p=9.5\times10^{-23}$). Under self-check every
evaluator remains significant after Holm adjustment; under steelman all five do
(Qwen weakest, Holm-adjusted $p=0.025$). Prompts and decoding
are identical across self- and cross-evaluation and never disclose authorship, so the
difference is not a framing artifact; self-preference through familiarity with one's
own phrasing remains a candidate mechanism for the excess. A second is
self-consistency: the diagonal evaluator may be re-deriving the same greedy solution
rather than recognizing authorship. Two families share a provider, so the matrix
measures model-family effects, not provider independence; adding a same-provider
indicator preserves the same-family effect (self-check odds ratio $2.06$,
$p=3.5\times10^{-10}$; steelman odds ratio $1.49$, $p=3.3\times10^{-6}$).
Cross rates also track error blatancy: the incorrect candidates of the highest-accuracy
generator (GLM, whose error stratum is smallest at $n_{T=0}=299$) are rejected more by
every evaluator than that evaluator's average over the other families (GLM self-check
cross false acceptance $73.2$--$82.5\%$). The blind spot is therefore mostly an
evaluator property with a measurable same-family excess on top.

\begin{table}[!t]
\centering
\caption{Self-versus-external evaluator discordance on matched incorrect candidates
(self-check arm). Rows are external evaluators, pooled over the four generator
families. ``Ext rej, own acc'' counts candidates the external evaluator rejects
while the generating family's own evaluator accepts (own family more permissive); the
reverse column counts the opposite direction. $p$ is the exact two-sided McNemar value;
the last column is the 95\% task-clustered paired-bootstrap interval for the net
discordance count. Pooled McNemar $p$-values are reported for reference; the bootstrap
is the primary inference.}
\label{tab:cross}
\resizebox{\columnwidth}{!}{%
\begin{tabular}{lcccc}
\hline
Evaluator & Ext rej, own acc & Ext acc, own rej & $p$ & Task-clustered 95\% CI \\
\hline
GPT-5.4 & 201 & 50 & $1.2\times10^{-22}$ & $[122,178]$ \\
DeepSeek-V4.1-Flash & 189 & 53 & $4.3\times10^{-19}$ & $[107,165]$ \\
GLM-5.3-Flash & 340 & 46 & $1.6\times10^{-56}$ & $[257,331]$ \\
Qwen-3.8-Flash & 148 & 59 & $5.2\times10^{-10}$ & $[63,116]$ \\
GPT-6-Luna & 199 & 53 & $4.5\times10^{-21}$ & $[114,175]$ \\
\hline
\end{tabular}}
\end{table}

\begin{table}[!t]
\centering
\caption{Full cross-evaluation false-acceptance matrix on incorrect candidates
(self-check / steelman, percent). Rows are evaluators, columns are generators;
diagonal cells are same-family evaluations. Denominators are arm-specific per cell
and are archived in \texttt{crossed\_model\_results.json}.}
\label{tab:cross_matrix}
\resizebox{\columnwidth}{!}{%
\begin{tabular}{lccccc}
\hline
Evaluator & GPT-5.4 & DeepSeek & GLM & Qwen & GPT-6-Luna \\
\hline
GPT-5.4 & \textbf{90.1/88.5} & 88.6/85.3 & 74.7/76.9 & 81.9/84.3 & 88.5/89.8 \\
DeepSeek & 79.0/79.4 & \textbf{96.9/94.4} & 80.4/77.2 & 87.1/89.2 & 86.7/87.4 \\
GLM & 72.8/73.3 & 78.8/77.4 & \textbf{88.8/85.5} & 77.2/74.7 & 72.7/73.1 \\
Qwen & 84.8/84.1 & 94.2/92.4 & 82.5/81.2 & \textbf{94.9/90.6} & 85.6/88.0 \\
GPT-6-Luna & 82.1/75.8 & 88.9/84.6 & 73.2/69.2 & 87.5/81.5 & \textbf{93.5/86.9} \\
\hline
\end{tabular}}
\end{table}

\subsection{Boundaries of the Measurement}
\label{sec:threats}

\paragraph{HumanEval error stratum.} Nine incorrect candidates make every error flip
worth 11 points of false-acceptance rate, and the paired bootstrap interval for the
non-oracle contrast includes zero. HumanEval supplies the clean paired demonstration of
the diagnostic ordering; all powered claims rest on MBPP.

\paragraph{Missing-record sensitivity.} The paired MBPP analysis uses 4,655 of 4,735
pairs; every excluded pair has exactly one ledger-recorded cause
(Appendix~\ref{app:accounting}). GPT-6-Luna's adjusted result is unchanged under every
extreme imputation of the missing pairs; Qwen's remains significant unless all three of
its missing $T{=}0$ pairs are simultaneously maximally adversarial. The full bounds are
in Appendix~\ref{app:accounting}.

\paragraph{HumanEval call variability.} Run-to-run verdict stability was not estimated
on HumanEval; the $4/9$ versus $3/9$ context-comparison contrast may partly reflect
model-call variability even at temperature zero. The fresh MBPP ledger measures verdict
stability directly (Section~\ref{sec:empirical}).

\paragraph{MBPP specification and interface mismatch.} Generators and evaluators saw
only the natural-language description and the entry-point name; the reference
asserts were never shown to either. The ground-truth label is therefore the benchmark's
correctness criterion, not a judgment about reasonable alternative readings of the
prose. The corpus is the frozen 947-task subset inherited from the earlier preserved
build (of the canonical 974; the 27 excluded task ids and their provenance are archived
in \path{data/MBPP.corpus-provenance.json}), and reference-assert counts per kept task are three in 936 of
the 947 tasks, with the remaining eleven carrying four to six (per-task reference
asserts are archived in \path{data/MBPP.jsonl}). The full-population severity audit shows 84.4\% of incorrect
candidates fail \emph{every} reference assert: the dominant pattern is wholesale
behavioral failure under the benchmark's own criterion, not a near-miss on one
ambiguous test. Six of the 1,987 (0.3\%) warrant an explicit reconciliation note:
the production harness applies a conservative static guard that rejects, before
execution, any candidate whose code references \texttt{compile()},
\texttt{getattr()}, or \texttt{import sys}, and these six were blocked at that
stage. Executed outside the guard, each passes every reference assert for its
task. Ground truth is defined by the operational harness, so the frozen labels
stand; but because the guard-free reading would move these six to the correct
pool, no pooled or per-family incorrect-candidate rate in this paper shifts by
more than $0.6$ percentage points under either labeling (the within-crash
TypeError share, a ratio of two moved counts, shifts by at most $0.7$ points).
The moderator analysis
has exactly one boundary stratum exposed to this relabeling: its Qwen code-length-high
stratum reports $68\%$ joint acceptance ($n=22$) under the frozen labels and would
read $66.7\%$ ($n=21$) under the guard-free reading; every other stratum with at
least ten paired incorrect candidates is unchanged in its relation to the floor. $\kappa$ is measured under that criterion throughout. A deterministic failure-type
breakdown of the 1,987 incorrect candidates finds only $0.2\%$ with a missing or
renamed entry point; $62.8\%$ run and fail on values, and $37.0\%$ land in the crash-or-guard-block
bucket (\path{data/mbpp_v2/analysis/failure_type_breakdown.json}). Of the 736 in that
bucket, $724$ are genuine runtime crashes and $12$ are pre-execution AST-guard blocks;
$81.5\%$ of the genuine crashes are \texttt{TypeError} (590/724;
\path{data/mbpp_v2/analysis/crash_exception_breakdown.json}). Missing or renamed entry points are therefore negligible, but this
breakdown does not distinguish genuine wrong answers from defensible alternative
readings of underspecified prose; the one-assert-visible Qwen control settles that
question directly: on all 947 tasks, generation accuracy is $56.1\%$ with one example
assert shown to both generator and evaluator, versus $55.9\%$ on the same tasks in
production, while self-check false acceptance is $92.5\%$ versus $94.9\%$ and
steelman false acceptance is $92.9\%$ versus $90.6\%$. Hidden conventions therefore
do not rescue generation or explain the evaluators' acceptance rates
(\path{data/mbpp_v2/assert_visible_qwen/assert_visible_analysis_summary.txt}).

\paragraph{Token-identical replicas.} The replication arm's two greedy candidates are
token-identical in $33.9\%$ of cells (up to $82\%$ for GPT-5.4); for those cells,
$T_B=T_A$ trivially and verdict repetition is partly guaranteed by construction. The
replication is therefore interpreted through the conditional-on-$T_A{=}0$ contrast and
the always-accept baseline, not through the headline Fisher associations.

\paragraph{Open by construction.} Neither experiment identifies $Z$ or measures the
latent information ceiling; that is the open problem stated in
Section~\ref{sec:falsifiable_z}, not a defect the data could have repaired. The
measured contrasts are properties of this two-arm battery under these prompts and the
strict verdict grammar.

\section{Design Principles for External Selection}
\label{sec:architecture}

We call evaluation \emph{external} when it uses a channel not fully controlled by the
generator---a formal tool, an independent model, or a fresh instance of the same model;
``external'' describes the channel, not the model. The contrapositive of the negative
results then motivates a design heuristic, offered as a design target rather than a
proved corollary:

\begin{quote}
\textbf{External Selection Design Heuristic.}
Evaluation is more likely to provide substantial information about correctness when
(1)~the selector accesses information not contained in $(G,Z)$, breaking the
conditional independence $S\perp T\mid(G,Z)$; (2)~the selector's blind spots have low
overlap with the generator's; and (3)~the selector satisfies the conditions of
Proposition~\ref{prop:multi_agent}: bounded false acceptance and true acceptance
exceeding that bound.
\end{quote}

\subsection{Selection Channels, Ranked by Independence Strength}

\paragraph{Formal verification.} Proof assistants (Lean, Coq, Isabelle) and type
checkers select under mathematical ground truth. The Prover Agent
framework~\cite{baba2025prover} illustrates the channel concretely: an informal
reasoning LLM coordinated with Lean achieved 88.1\% accuracy on MiniF2F, a new
state of the art among methods using small language models at a substantially lower
sample budget~\cite{baba2025prover}.
Lean applies a criterion external to the language model's own confidence or reasoning
trace---a qualitatively stronger source of independence than same-model context
separation.

\paragraph{Executable verification.} Unit tests, property-based tests, and simulation
checks select under computational ground truth; the interpreter does not care how
confident the generator was. This is the channel instantiated by the execution-grounded
control in Section~\ref{sec:empirical}; it shows zero observed false acceptances there,
although the verdicts remain model outputs (it also falsely rejects six PASS
candidates); zero false acceptances in nine cases is an observed result, not a
guarantee, and the Wilson 95% interval $[0.000,0.299]$ leaves the underlying rate
substantially unconstrained. It marks what a
direct computational correctness channel buys.

\paragraph{Numerical invariants and retrieval grounding.} Dimensional analysis,
conservation laws, and sanity bounds select under physical ground truth; citation and
quote verification against a fixed corpus selects under documentary ground truth.

\paragraph{Independent critics.} Models with different training data, architectures, or
optimization objectives have partially independent failure modes. The cross-evaluation
matrix of Section~\ref{sec:mechanism} measures exactly this channel: external
evaluators reject errors the generating family accepted at $2.5$--$7.4$ to one,
but their absolute false-acceptance rates remain high, so independence is partial, not
complete.

\paragraph{Fresh-context evaluation as an engineering heuristic.} The same model under
fresh context removes the generation trace but not the shared weights and training
distribution, so it cannot establish the independence conditions above. The HumanEval
context comparison provides no compelling support for a large fresh-context benefit
and is inconclusive about moderate effects; \cite{tsui2025selfcorrection}'s ``Wait''
intervention reduced the blind spot by 89.3\% through context perturbation, a
performance result whose mechanism is not established. Context separation remains a
low-cost intervention that alters the activation context even when it cannot alter the
underlying weights, and is recommended on that basis rather than as a validated remedy.

\begin{Definition}[Context separation]
\label{def:context_sep}
Two evaluation contexts are \emph{separated} if the evaluating context has no access to:
(1) the generation trace (intermediate reasoning steps), (2) the prompt scaffolding
(instructions that shaped generation), or (3) hidden state from the generation process.
\end{Definition}

\subsection{Why Independence Matters}

Perfect independence is not required, but reduced joint false acceptance alone is not
sufficient for evidence compounding. For unanimous acceptance, define
$\tau_{1:k}=\Pr(A_{1:k}=1\mid T=1)$, $\kappa_{1:k}=\Pr(A_{1:k}=1\mid T=0)$, and
$\Lambda_{1:k}=\tau_{1:k}/\kappa_{1:k}$. For positive rates,
\begin{equation}
\frac{\Lambda_{1:k}}{\Lambda_{1:k-1}}=
\frac{\Pr(A_k=1\mid A_{1:k-1}=1,T=1)}
{\Pr(A_k=1\mid A_{1:k-1}=1,T=0)},
\end{equation}
so the $k$-th acceptance adds positive incremental evidence exactly when this ratio
exceeds one: joint true acceptance must decrease by a smaller factor than joint false
acceptance. If decisions are mutually conditionally independent given $T$, then
$\Lambda_{1:k}=\prod_{i=1}^k(\tau_i/\kappa_i)$, and a selector with $\tau_k>\kappa_k$
adds positive evidence. Without that factorization, the conditional ratio above is the
criterion: individual rates $(\tau_1,\kappa_1)=(0.9,0.6)$ give $\Lambda_1=1.5$, whereas
joint rates $(0.2,0.3)$ give $\Lambda_{1:2}=2/3$ despite reduced joint false
acceptance. Diverse multi-agent panels can therefore help when their aggregation
improves these likelihood ratios~\cite{du2023debate, liang2023debate}; reduced joint
failure probability alone does not establish that mechanism. Persona-based
diversity---casting the evaluator as a rigorous mathematician, a skeptical physicist,
or a hostile referee---is a cheap single-model attempt at the same effect, since each
persona predicts a different distribution of expert behavior with different blind
spots; whether it yields acceptance rates satisfying
Proposition~\ref{prop:multi_agent} is an open question.

\subsection{Architecture and Mitigations}

The principles organize around three functions: high-entropy generation to explore the
hypothesis space, external selection to filter candidates under criteria the
generator's blind spots cannot corrupt, and feedback to update generation based on what
survives selection. The selector's criteria form a hierarchy: formal (does it compile,
prove, or type-check?), executable (does it pass tests?), numerical (does it satisfy
invariants?), grounded (do citations check out?), and adversarial (does it survive
independent critique?). The feedback component records which constraints killed
candidates, preserves surviving patterns as templates, and escalates ambiguous
survivors to human review. This is a design proposal, not a validated system:
Section~\ref{sec:empirical} tests the evaluation components, not the integrated
architecture.

The failure modes of such a pipeline are themselves selection problems. Correlated
critics reintroduce the original blind spot through the back door; prompt injection on
the selection channel can manipulate verdicts; spoofable external checks (doctored
test harnesses, retrieval poisoning) convert an independence assumption into an attack
surface; and feedback gaming lets generators optimize for the selector rather than for
correctness. The mitigations follow the same independence principle: diverse selection
criteria with periodic architecture audits, independent sanitization before evaluation,
multiple independent checkers, human review of accepted candidates, and criteria
rotation so the feedback channel cannot be overfit.

\section{Related Work}
\label{sec:related}

\paragraph{LLM self-correction.} \cite{huang2024selfcorrect} provide the foundational
negative result: LLMs struggle to self-correct without external feedback, and in some
settings self-correction degrades performance. \cite{tyen2024llmsstepbystep} find that
models often fail to locate logical errors in reasoning, even in objective settings,
while their ability to correct an error improves substantially when the error location
is supplied. \cite{tsui2025selfcorrection}
supply the closest empirical anchor: models correct identical errors presented as
external input but fail on the same errors in their own outputs, at an average 64.5\%
failure rate across 14 models on three purpose-built benchmarks (SCLI5, GSM8K-SC,
PRM800K-SC), and the ``Wait'' intervention reduced the blind spot by 89.3\%. We treat
these findings as empirical context rather than formal validation: the mapping between
``failure to self-correct'' and the theorem's formal quantities is not established, and
the mechanism by which ``Wait'' produces its effect is not identified by the
performance evidence. Where we provide theoretical conditions under which
self-evaluation may fail, Tsui provides empirical measurement of failure rates and a
candidate intervention whose mechanism remains to be characterized.

\paragraph{Correlated errors and evaluator bias.} \cite{kim2025correlated} find,
across 350+ LLMs, that model pairs agree on the same wrong answer approximately 60\%
of the time when both err, with shared provider, shared architecture, and higher
capability each associated with higher error correlation, alongside larger and more
accurate models. Their cross-generator
agreement measurement is not equivalent to an evaluator accepting a generator's
incorrect answer, so it does not directly instantiate our high-$\kappa$ condition; it
motivates checking for coupled evaluation failures. \cite{wataoka2024selfpreference}
and \cite{li2024llmjudgebias} document self-preference and related LLM-as-a-judge
biases, consistent with the self-leniency excess measured in our cross-evaluation
matrix.

\paragraph{Verification and external feedback.} \cite{wang2023selfconsistency}
demonstrated that sampling multiple reasoning paths and taking the majority answer
improves accuracy. The standard Condorcet-style calculation gives sufficient conditions
for aggregation gains under its independent-error and voter-competence assumptions, not
necessary conditions for every joint voting model: correlated samples can still improve
a vote, and a scalar error correlation does not by itself determine the majority-error
probability. If every sample deterministically repeats the same answer, voting adds no
diversity. More generally, the effect depends on the joint distribution of answers and
errors, including whether incorrect votes concentrate on one alternative. Our analysis
therefore motivates measuring that joint failure structure rather than assuming a
universal monotone gain from decorrelation. Multi-agent debate and verification
frameworks~\cite{du2023debate, liang2023debate} introduce altered evaluation channels,
with effectiveness depending on how much independent information they supply;
\cite{gou2023critic}'s CRITIC framework showed that tool augmentation improves
evaluation through external feedback, and \cite{chen2024reconcile} found that model
diversity among agents improves reconciliation.
\cite{lightman2023process} showed that
per-step process supervision outperforms outcome supervision in their
mathematical-reasoning setting: an example of an external signal entering at a finer
granularity than final-answer self-evaluation.
\cite{krogh1995ensemble} established the classical ensemble-diversity result on which
the independence argument rests, and \cite{cover2006elements} supply the
information-theoretic background. \cite{alphafold} provides an analogy from a domain in which predictions are evaluated
against experimentally determined protein structures and quantitative structural
benchmarks; it does not itself test the self-evaluation mechanism studied here.

\paragraph{Apparent counterexamples.} Successful iterative-improvement systems are
not necessarily counterexamples to this diagnostic. Self-Refine~\cite{madaan2023selfrefine}
studies iterative self-feedback across generation and reasoning tasks;
Reflexion~\cite{shinn2023reflexion} augments agents with verbal feedback and memory;
and Constitutional AI~\cite{bai2022constitutional} uses an explicit constitution and
AI-generated critique and preference signals. These settings differ from the narrower
question studied here: whether an evaluator can reliably assign a correctness verdict
to a frozen candidate without an independently grounded correctness channel. When
\cite{huang2024selfcorrect} removed oracle feedback from self-correction setups,
improvements vanished.

\section{Conclusions}

The paper separates three questions that are often conflated: whether a model can
generate a correct answer, whether it can recognize its own errors under an internal
evaluation procedure, and how its selection behavior changes when an external
correctness channel is introduced. On the five MBPP
families, internal evaluation accepts most incorrect candidates even after structured
critique, and accepted candidates remain only $55$--$70\%$ correct (on the smaller
HumanEval pilot, steelman accepts $3/9$ incorrect but accepted candidates are
$96$--$98\%$ correct); only the external correctness channel sharply separates correct
from incorrect. Steelman evaluation
reduces false acceptance in every family, but the reduction is protocol-specific and
mechanistically split: Qwen's neutral re-examination prompt fails to reproduce the
steelman reduction at matched token use, while GPT-6-Luna's neutral and steelman rates
are nearly equal (not an equivalence result). The cross-evaluation matrix shows the permissiveness extends to other
families' incorrect candidates, with a same-family
acceptance excess in every evaluator under self-check and in all five under
steelman. Repeated unanimous acceptance carries $27$--$44\%$ residual
error at $k{=}3$, and joint acceptance between the arms is largely what the marginal
rates alone predict.

The theory is conditional and concerns average information: under joint conditional
independence, $k$ rounds of self-critique provide mutual information bounded by what
the shared latent variable $Z$ mediates; the separate miscalibration proposition shows
how an observer treating coupled evaluations as independent can become arbitrarily
confident in errors; and the positive-evidence proposition gives the constructive counterpart,
a quantifiable evidence floor from measurable acceptance rates. Neither experiment
identifies $Z$ or measures the latent ceiling, and constructing an independently
specified, informative $Z$ remains the key open problem. The framework's value is
diagnostic rather than absolute: it identifies the failure quantity that matters---acceptance
of incorrect outputs---and the design principle that follows: validate through a channel with independent
correctness access.

\paragraph{Open problems.} Constructing an independently specified, informative $Z$
(Section~\ref{sec:falsifiable_z}); a pre-specified self-defense control alongside the
steelman, with realized tokens and cost reported and the primary paired comparison
pre-registered; testing the joint conditional-independence relation under explicit
candidate-level coarsening; verdict-stability estimation at small scale; and measuring
whether persona-based diversity yields acceptance rates satisfying
Proposition~\ref{prop:multi_agent}.

\section*{Acknowledgment}
The author thanks the anonymous reviewers for helpful feedback and Joe Reid for
infrastructure buildout and design. This research received no external funding. The
empirical measurements used the public HumanEval and MBPP benchmarks. The preserved
benchmark harness, ledgers, and audit artifacts are publicly available at
\texttt{https://github.com/\allowbreak AndBrilliant/\allowbreak selfcorr-benchmarks}
(ledgers frozen since commit \texttt{81537cf}; audited release tagged \texttt{v1.0-audited}). The author declares no conflicts of interest.

Use of AI tools: AI-assisted tools were used during manuscript preparation for language
editing and drafting support. All ideas, methodology, theorems, and technical content
are the author's own work, and the author takes full responsibility for the accuracy,
originality, and integrity of the submitted content.

\appendices

\section{Ledger Accounting and Sensitivity}
\label{app:accounting}

The paired MBPP analysis uses 4,655 of 4,735 pairs; every excluded pair is accounted
for by a single cause recorded in the ledger (Table~\ref{tab:accounting}). The 67
output-ceiling pairs (17 generation-side, 50 evaluator-side) are restored malformed
records: their final attempt terminated at the provider's output ceiling
(\texttt{finish\_reason=length}), a deterministic degenerate output pattern at
temperature zero, and they were never retried or recoded. The 80 pairs excluded from
the paired analysis additionally include 13 pairs whose API call completed cleanly but
whose response did not match the strict verdict grammar (9 GPT-5.4 steelman, 4
GPT-6-Luna self-check): extraction failures, not provider failures, kept visible under
the same rule. Arm-level, the ledger's 9,470 evaluator arms are 9,370 parsed, 66
malformed, and 34 not-run (the 17 generation-failed pairs' two arms each). No record
is dropped silently in any category. Two supplementary ledgers extend the battery
under the same protocol: the neutral-verification control arm (4,718 records, 6,965
events, 29 malformed) and the cross-evaluation matrix (18,872 records, 38,295 events,
399 malformed), each with its own SHA-256 manifest. During the control run the OpenAI
key was invalidated mid-run; the 747 affected GPT-5.4 tasks were re-run after key
replacement, and the auth-failure records remain preserved in
\path{auth_failed_records.jsonl} alongside the event-id repair log.

\begin{table}[!t]
\centering
\caption{MBPP fixed-denominator accounting, 4,735 pairs. Categories are
disjoint and exhaustive; every excluded pair has exactly one ledger-recorded
cause. ``Grammar'' means the API call completed cleanly but the response
carried no strict terminal verdict.}
\label{tab:accounting}
\resizebox{\columnwidth}{!}{%
\begin{tabular}{lrrrrr}
\hline
Family & Paired & Generation & \multicolumn{3}{c}{Evaluator-arm exclusions} \\
 & & (ceiling) & ceiling & grammar & arms \\
\hline
GPT-5.4 & 938 & 0 & 0 & 9 & sc:0, st:9 \\
DeepSeek-V4.1-Flash & 917 & 9 & 21 & 0 & sc:9, st:10, both:2 \\
GLM-5.3-Flash & 917 & 5 & 25 & 0 & sc:15, st:10 \\
Qwen-3.8-Flash & 940 & 3 & 4 & 0 & st:3, both:1 \\
GPT-6-Luna & 943 & 0 & 0 & 4 & sc:4 \\
\hline
Total & 4,655 & 17 & 50 & 13 & 80 excl.\\
\hline
\end{tabular}}
\end{table}

Because missing verdicts are excluded rather than imputed, the paired FAR McNemar
contrast is bounded under the most adversarial extreme imputation (every missing pair
assigned to maximize discordance against steelman). GPT-6-Luna's adjusted result is
unchanged under any imputation of its missing pairs (worst-case
$p_{\mathrm{Holm}}=3.7\times10^{-5}$). Qwen's survives unless all three of its missing
$T{=}0$ pairs are simultaneously maximally adversarial (worst-case
$p_{\mathrm{Holm}}=0.080$; any imputation with at most two adversarial pairs keeps
$p_{\mathrm{Holm}}\le0.045$). The three non-significant families remain
non-significant under every imputation.

\section{A Predictive Interpretation}
\label{app:mechanism}

The information-theoretic results establish \emph{that} self-evaluation can fail under
correlated error; this appendix offers one interpretation of \emph{why} such
correlation may be elevated in language models. The account is consistent with the
formal analysis but does not depend on it.

\cite{bender2021parrots} characterize language models as systems that stitch together
sequences of linguistic forms according to probabilistic information about how they
combine, without reference to meaning. Under this framing, when asked to evaluate a
hypothesis, a language model predicts what evaluative text would likely follow the
prompt, inheriting whatever patterns characterize human evaluation behavior in the
training distribution: prestige deference, format heuristics, social smoothing, and
narrative continuation. Alignment training shifts \emph{which} human behavior is
predicted but may not change the underlying operation: \cite{sharma2024sycophancy}
demonstrate that RLHF-trained models systematically exhibit sycophancy and that human
preference judgments favor sycophantic responses, creating a training signal toward
agreement. Commercial systems commonly use ``helpful assistant'' role framing rather than a
rigorous evaluation protocol~\cite{zheng2023helpful}.

One observable phenomenon follows: \emph{format consistency}. We hypothesize that
submitting manuscripts of widely varying quality to major language models with neutral
prompts will yield near-uniform response format---balanced positive and negative points
regardless of input quality. This hypothesis is stated in falsifiable form to invite
experimental follow-up; it is not an established result. The behaviors above are
defaults, not constraints: prompt design substantially affects model behavior, with
performance differences of up to 76 percentage points in the few-shot LLaMA-2-13B
experiments of \cite{sclar2024formatting} and accuracy gains from role
framing~\cite{kong2023roleplay}. If generation runs under a ``helpful assistant''
framing and evaluation under the same framing, error-correlation risk is structurally
elevated, since both are shaped by the same training signal; context separation may
help because a fresh context with explicit critic framing resets the prediction
target, potentially reducing the shared context-level component of $Z$.


\section{Prompts and Decoding Settings}
\label{app:prompts}

All MBPP calls used temperature zero. The request payloads contained no
reasoning-effort field and no output-length cap; provider defaults therefore apply.
Model identifiers were \texttt{gpt-5.4} (the API alias resolving to the pinned
\texttt{gpt-5.4-2026-03-05} snapshot),
\texttt{deepseek/deepseek-v4.1-flash}, \texttt{z-ai/glm-5.3-flash},
\texttt{qwen/qwen3.8-flash}, and \texttt{openai/gpt-6-luna}. The HumanEval
snapshot and reasoning setting are given in Section~\ref{sec:empirical}.

The MBPP self-check prompt was:
\begin{quote}\scriptsize\ttfamily\raggedright
You are evaluating one immutable Python candidate solution.

[PROBLEM]
{task prose and entry-point name}

[CANDIDATE]
{frozen code}

Analyze whether the candidate satisfies the specification. Do not modify,
repair, rewrite, or provide an alternative solution. End with exactly
one final line containing only one of:
VERDICT: ACCEPT
VERDICT: REJECT
\end{quote}

The MBPP steelman prompt was:
\begin{quote}\scriptsize\ttfamily\raggedright
You are a neutral arbiter judging a debate between two expert Python programmers.

[PROBLEM]
{task prose and entry-point name}

[CANDIDATE]
{frozen code}

Do the following in order:
1. STEELMAN AGAINST: give the strongest case that the candidate is wrong.
2. STEELMAN FOR: give the strongest case that the candidate is correct.
3. Decide which case survives scrutiny. Judge only the candidate as given;
   do not provide a corrected version.
End with exactly one final line containing only one of:
VERDICT: ACCEPT
VERDICT: REJECT
\end{quote}

The neutral-control prompt was:
\begin{quote}\scriptsize\ttfamily\raggedright
You are performing a careful neutral verification of one immutable Python
candidate solution.

[PROBLEM]
{task prose and entry-point name}

[CANDIDATE]
{frozen code}

Verify the candidate against the specification step by step: check the logic,
trace the behavior on the requirements, and confirm whether it satisfies the
specification. Do not modify, repair, rewrite, or provide an alternative
solution, and do not presume the candidate is correct or incorrect; verify it.
End with exactly one final line containing only one of:
VERDICT: ACCEPT
VERDICT: REJECT
\end{quote}

The production self-check wording said ``specification and tests''; the reference
asserts were nevertheless not shown to the evaluator, so ``tests''
referred to the hidden ground-truth criterion rather than visible evidence. The
one-assert-visible control adds one example assert to both generator and
evaluator. The cross-evaluation arms used the same self-check and steelman
prompts verbatim for every evaluator--generator pair.

\bibliographystyle{BibTeXtran}
\bibliography{refs}

\end{document}

